\documentclass[11pt]{article}

\usepackage[T1]{fontenc}
\usepackage[margin=0.82in]{geometry}
\usepackage{amsmath,amssymb}
\usepackage{booktabs}
\usepackage{array}
\usepackage{xcolor}
\usepackage{hyperref}
\usepackage{fancyhdr}
\usepackage{enumitem}

\definecolor{navy}{HTML}{17365D}
\definecolor{blue}{HTML}{2F75B5}
\definecolor{pale}{HTML}{EAF2F8}
\definecolor{green}{HTML}{2E7D32}
\definecolor{gray}{HTML}{555555}

\hypersetup{
  colorlinks=true,
  linkcolor=navy,
  urlcolor=blue,
  pdftitle={Nonlinear-Manifold ROM: Poisson and Burgers from Training to Prediction},
  pdfauthor={Tunable NM-ROM project}
}

\pagestyle{fancy}
\fancyhf{}
\lhead{NM-ROM: Poisson and Burgers}
\rhead{Training $\rightarrow$ precompute $\rightarrow$ solve}
\lfoot{\scriptsize $P$: grid values \quad $K$: latent variables \quad $R$: learned patterns \quad $M$: test equations}
\rfoot{\thepage}
\setlength{\headheight}{14pt}

\setlength{\parindent}{0pt}
\setlength{\parskip}{0.55em}
\setlist[itemize]{itemsep=0.15em,topsep=0.25em}
\setlist[enumerate]{itemsep=0.2em,topsep=0.3em}

\newcommand{\mat}[1]{\mathbf{#1}}
\newcommand{\vect}[1]{\mathbf{#1}}
\newcommand{\R}{\mathbb{R}}
\newcommand{\keyidea}[1]{%
  \begin{center}
  \fcolorbox{blue}{pale}{\begin{minipage}{0.91\textwidth}
  #1
  \end{minipage}}
  \end{center}}

\begin{document}

\begin{center}
  {\LARGE\bfseries\color{navy} Nonlinear-Manifold ROM (NM-ROM)}\\[0.35em]
  {\Large From Training to Poisson and Burgers Prediction}\\[0.8em]
  {\color{gray} A numerical walkthrough of the separable-decoder architecture}
\end{center}

\vspace{0.5em}

\keyidea{\textbf{Purpose of this document.}
Part I follows Poisson from training to a final prediction with small numbers.
Part II explains the same architecture as a general nonlinear-manifold ROM.
Part III shows exactly what changes for time-dependent Burgers.  The Poisson
project dimensions in Section~\ref{sec:real-sizes} and the Burgers dimensions
in Section~\ref{sec:burgers-sizes} come from the implemented configurations.
The small arithmetic examples are teaching numbers, not measured results.}

\setcounter{tocdepth}{1}
\tableofcontents
\newpage

\section{The problem and the model}

We want to solve the Poisson equation
\begin{equation}
  -\Delta u = f,
  \qquad
  u=0 \text{ on the boundary}.
\end{equation}
Here $f$ is a new source field and $u$ is the solution we want to predict.

The learned decoder has two separate neural networks:
\begin{equation}
  u(\vect{x};\vect{z})
  = s\,b_c(\vect{x})\,
    \vect{g}(\vect{x})^{\mathsf T}\vect{h}(\vect{z}).
\end{equation}

The two networks have different jobs:
\begin{itemize}
  \item $\vect{g}(\vect{x})$ produces $R$ spatial-pattern values at coordinate
        $\vect{x}$.  Here $R$ means the number of learned spatial patterns.
  \item $\vect{h}(\vect{z})$ produces $R$ coefficients, one coefficient for
        each of those $R$ learned spatial patterns.
  \item $b_c(\vect{x})$ is zero on the boundary, so every decoded field
        satisfies the boundary condition automatically.
\end{itemize}

On all interior grid points at once, the decoder becomes
\begin{equation}
  \vect{u}(\vect{z}) = \mat{G}\,\vect{h}(\vect{z}),
  \qquad
  \mat{G}\in\R^{P\times R},
\end{equation}
Here $P$ is the number of interior grid values in the solution vector, and
$R$ is the number of learned spatial patterns.  Thus $\mat{G}$ has one row per
grid value and one column per learned pattern.  In the implementation, row
$i$ of $\mat{G}$ is
\begin{equation}
  G_{i,:}
  =s\,b_c(\vect{x}_i)\,\vect{g}(\vect{x}_i)^{\mathsf T},
\end{equation}
so the fixed training-data scale $s$ and the hard boundary mask are already
inside the cached bank.

\keyidea{The new forcing $f$ is \emph{not} fed into $\vect{h}$.
For a new problem, $f$ enters the residual.  A numerical solver changes
$\vect{z}$ until the decoded field satisfies the projected Poisson equation.}

\keyidea{\textbf{Three sizes to remember throughout the document.}
$P$ counts physical-grid values in $\vect{u}$; $K$ counts the latent variables
in $\vect{z}$ that the online solver changes; and $R$ counts both the learned
spatial patterns in $\mat{G}$ and their coefficients in $\vect{h}(\vect{z})$.
The same reminder appears in every page footer.}

\section{The real project dimensions}
\label{sec:real-sizes}

For the $N=256$ Poisson quadrature-free run, the full mesh has
$256\times256$ points.  Removing the outer boundary leaves
\begin{equation}
  P=(256-2)^2=254^2=64{,}516
\end{equation}
interior points.  In other words, $P$ is the length of a full interior solution
vector, not the number of latent variables or learned patterns.

The main dimensions are:

\begin{center}
\begin{tabular}{@{}lll@{}}
\toprule
Object & Shape & Meaning \\
\midrule
$\vect{z}$ & $16=K$ & latent variables changed by the solver \\
$\vect{h}(\vect{z})$ & $64=R$ & one coefficient per learned spatial pattern \\
$\mat{G}$ & $64{,}516\times64=P\times R$ & grid values by learned patterns \\
$\mat{\Phi}^{\mathsf T}$ & $64\times64{,}516$ & 64 sine test functions \\
$\mat{B}$ & $64\times64$ & projected spatial bank \\
$\mat{A}$ & $64\times64$ & projected Poisson operator \\
$\vect{r}(\vect{z})$ & $64$ & residual seen by the latent solver \\
$\mat{J}_r$ & $64\times16$ & residual Jacobian with respect to $\vect{z}$ \\
\bottomrule
\end{tabular}
\end{center}

The training configuration used 512 training solutions and 100,000 optimizer
steps, with
\begin{equation}
  K=16,\qquad R=64,\qquad M=64.
\end{equation}
That means 16 latent variables are adjusted online ($K=16$), the decoder uses
64 learned spatial patterns and 64 matching coefficients ($R=64$), and the
residual retains 64 test equations ($M=64$).

\textbf{Notation warning.}  The blackboard-bold symbol $\mathbb{R}$ in an
expression such as $\mathbb{R}^{16}$ means ``vectors of real numbers.''  It is
not the model-size symbol $R$.  For example, $\mathbb{R}^{16}$ below is a
real-valued vector with 16 entries, while $R=64$ counts learned patterns.

The spatial path begins with 64 learned frequency vectors, expanded into 64
sine values and 64 cosine values.  Its architecture is
\begin{equation}
  \R^2
  \longrightarrow 128\text{ Fourier features}
  \longrightarrow 128
  \longrightarrow 128
  \longrightarrow \R^{64},
\end{equation}
with SiLU activations after the two hidden dense layers and a linear final
layer.  Its output is multiplied by the fixed training-data RMS scale and the
hard boundary mask when the bank is formed.  The latent path is
\begin{equation}
  \R^{16}\longrightarrow128\longrightarrow128\longrightarrow\R^{64},
\end{equation}
again with SiLU hidden activations and a linear final layer.  A separate linear
map from $\R^{16}$ to $\R^{64}$ is added as a skip connection.  In this path,
the 16 input entries are the $K=16$ latent variables, and the 64 output entries
are the $R=64$ learned-pattern coefficients.

\section{Offline training}

\subsection{Step 1: generate solution snapshots}

The full-order Poisson solver produces training pairs
\begin{equation}
  (\vect{f}_s,\vect{u}_s),
  \qquad s=1,\ldots,512.
\end{equation}
Each $\vect{u}_s$ contains $P=64{,}516$ interior solution values in the
$N=256$ run; here $P$ again means the number of physical-grid values stored in
one solution snapshot.

\subsection{Step 2: assign a trainable code to every snapshot}

Because this is an auto-decoder, there is no encoder.  Every training snapshot
gets its own trainable latent vector:
\begin{equation}
  \vect{z}_s\in\R^{16}.
\end{equation}
The 16 entries of this vector are the $K=16$ latent variables assigned to that
snapshot.  The 512 latent codes and the weights of both neural networks are
optimized together.

\subsection{Step 3: reconstruct every training solution}

At a training step, evaluate the spatial bank below.  It has $P$ rows, one for
each interior grid value, and $R$ columns, one for each learned spatial pattern:
\begin{equation}
  \mat{G}_{\theta_g}
  = \begin{bmatrix}
      s\,b_c(\vect{x}_1)\vect{g}_{\theta_g}(\vect{x}_1)^{\mathsf T}\\
      \vdots\\
      s\,b_c(\vect{x}_P)\vect{g}_{\theta_g}(\vect{x}_P)^{\mathsf T}
    \end{bmatrix}
  \in\R^{P\times R}
\end{equation}
The coefficient matrix below has 512 rows, one for each training snapshot, and
$R$ columns, one coefficient for each learned spatial pattern:
\begin{equation}
  \mat{H}_{\theta_h}
  = \begin{bmatrix}
      \vect{h}_{\theta_h}(\vect{z}_1)^{\mathsf T}\\
      \vdots\\
      \vect{h}_{\theta_h}(\vect{z}_{512})^{\mathsf T}
    \end{bmatrix}
  \in\R^{512\times R}.
\end{equation}

All training reconstructions are produced together by the following matrix.
It has 512 rows, one reconstructed snapshot per row, and $P$ columns, one for
each interior grid value:
\begin{equation}
  \widehat{\mat{U}}=\mat{H}_{\theta_h}\mat{G}_{\theta_g}^{\mathsf T}
  \in\R^{512\times P}.
\end{equation}

The main loss is the global relative mean-square reconstruction error,
\begin{equation}
  \mathcal{L}_{\mathrm{recon}}
  =
  \frac{\lVert\widehat{\mat{U}}-\mat{U}\rVert_F^2}
       {\lVert\mat{U}\rVert_F^2},
\end{equation}
plus a small feature-Gram regularizer that discourages badly conditioned or
redundant columns in $\mat{G}$.

Adam updates
\begin{equation}
  \theta_g,\qquad \theta_h,\qquad
  \vect{z}_1,\ldots,\vect{z}_{512}.
\end{equation}

\subsection{A tiny training calculation}

For the worked example, use only four interior points, two latent variables,
and three spatial features:
\begin{equation}
  P=4,\qquad K=2,\qquad R=3.
\end{equation}
Thus $P=4$ means four physical-grid values in each solution, $K=2$ means two
latent variables for the solver to adjust, and $R=3$ means three learned
spatial patterns and three matching coefficients.

Suppose one training target is
\begin{equation}
  \vect{u}_{\mathrm{train}}
  =\begin{bmatrix}1&0&1&2\end{bmatrix}^{\mathsf T}.
\end{equation}
An early prediction might be
\begin{equation}
  \widehat{\vect{u}}
  =\begin{bmatrix}0.3&-0.1&0.2&0.4\end{bmatrix}^{\mathsf T}.
\end{equation}
Its pointwise error and plain mean-square error are
\begin{align}
  \widehat{\vect{u}}-\vect{u}_{\mathrm{train}}
    &=\begin{bmatrix}-0.7&-0.1&-0.8&-1.6\end{bmatrix}^{\mathsf T},\\
  \mathrm{MSE}
    &=\frac{0.49+0.01+0.64+2.56}{4}=0.925.
\end{align}
This plain MSE is only a small arithmetic illustration; the production loss
uses the relative normalization written above.

After many updates, suppose training has produced the bank
\begin{equation}
\label{eq:Gtoy}
  \mat{G}
  =\begin{bmatrix}
      1&0& 1\\
      0&1& 1\\
      1&1& 0\\
      2&1&-1
    \end{bmatrix}.
\end{equation}
For the training snapshot above, suppose the trained head gives
\begin{equation}
  \vect{h}(\vect{z}_{\mathrm{train}})
  =\begin{bmatrix}1&0&0\end{bmatrix}^{\mathsf T}.
\end{equation}
Then
\begin{equation}
  \mat{G}\vect{h}(\vect{z}_{\mathrm{train}})
  =\begin{bmatrix}1&0&1&2\end{bmatrix}^{\mathsf T},
\end{equation}
so this teaching snapshot is reconstructed exactly.

\subsection{Step 4: freeze the learned model}

After training, the spatial network, Fourier frequencies, latent network, and
linear skip are frozen.  The training codes are no longer optimized.  They
may still be useful for constructing an initial guess for a new solve.

\section{One-time Poisson precomputation}

\subsection{Step 5: evaluate the frozen spatial bank}

Evaluate the frozen spatial network at every interior point, multiply by the
fixed output scale and hard boundary mask, and build a matrix with $P$ rows
(one per interior grid value) and $R$ columns (one per learned spatial pattern):
\begin{equation}
  \mat{G}\in\R^{P\times R}.
\end{equation}
This is the large spatial calculation.  It is not repeated during each latent
solver iteration.

\subsection{Step 6: construct the sine test functions}

Let
\begin{equation}
  \mat{\Phi}
  =\begin{bmatrix}\vect{\phi}_1&\cdots&\vect{\phi}_M\end{bmatrix}
  \in\R^{P\times M}
\end{equation}
contain the selected discrete sine modes.  The $P$ rows place every test
function on all $P$ interior grid values; the $M$ columns are the $M$ retained
test functions.  They are eigenvectors of the
discrete negative Laplacian $\mat{L}$:
\begin{equation}
  \mat{L}\vect{\phi}_m=\lambda_m\vect{\phi}_m.
\end{equation}

For a general fixed linear operator, the projected matrix would be built as
\begin{equation}
  \mat{A}=\mat{\Phi}^{\mathsf T}\mat{L}\mat{G}.
\end{equation}
For Poisson sine modes, symmetry and the eigenvector identity give the
shortcut
\begin{equation}
  \mat{A}=\mat{\Lambda}\mat{B},
  \qquad
  \mat{B}=\mat{\Phi}^{\mathsf T}\mat{G},
\end{equation}
where
\begin{equation}
  \mat{\Lambda}=\operatorname{diag}(\lambda_1,\ldots,\lambda_M).
\end{equation}

\subsection{Step 7: compute every small matrix in the toy example}

Take a $4\times4$ full grid on $[0,1]^2$.  It has a $2\times2$ interior, so
$P=4$ physical interior-grid values, and its grid spacing is $\Delta x=1/3$.
For readable arithmetic,
manually select two test modes, $M=2$.  Their values are
\begin{equation}
  \mat{\Phi}^{\mathsf T}
  =\begin{bmatrix}
      0.5& 0.5&0.5& 0.5\\
      0.5&-0.5&0.5&-0.5
    \end{bmatrix}.
\end{equation}
These are actual normalized two-dimensional discrete sine modes on this tiny
grid.  The production mode selector retains complete eigenvalue shells.  On
this unusually small grid, requesting two modes would therefore retain the
second mode's tied partner as a third mode.  The teaching example intentionally
omits that tied partner to keep every displayed matrix two rows tall; each
displayed vector remains a valid sine eigenmode.  Their negative-Laplacian
eigenvalues are
\begin{equation}
  \lambda_1=18,
  \qquad
  \lambda_2=36,
  \qquad
  \mat{\Lambda}=\begin{bmatrix}18&0\\0&36\end{bmatrix}.
\end{equation}

Projecting the bank in Equation~\eqref{eq:Gtoy} gives
\begin{align}
  \mat{B}
  &=\mat{\Phi}^{\mathsf T}\mat{G}\\[0.25em]
  &=
  \begin{bmatrix}
      0.5& 0.5&0.5& 0.5\\
      0.5&-0.5&0.5&-0.5
    \end{bmatrix}
  \begin{bmatrix}
      1&0& 1\\
      0&1& 1\\
      1&1& 0\\
      2&1&-1
    \end{bmatrix}\\[0.25em]
  &=\begin{bmatrix}
      2& 1.5&0.5\\
      0&-0.5&0.5
    \end{bmatrix}.
\end{align}

For example, the upper-left entry is
\begin{equation}
  B_{11}=0.5(1)+0.5(0)+0.5(1)+0.5(2)=2.
\end{equation}

The projected Poisson matrix is
\begin{align}
  \mat{A}
  &=\mat{\Lambda}\mat{B}\\
  &=\begin{bmatrix}18&0\\0&36\end{bmatrix}
    \begin{bmatrix}2&1.5&0.5\\0&-0.5&0.5\end{bmatrix}\\
  &=\begin{bmatrix}
      36& 27& 9\\
       0&-18&18
    \end{bmatrix}.
\end{align}

For the project choice $\alpha=1$, the residual weighting is the inverse
eigenvalue scaling
\begin{equation}
  \mat{W}=\mat{\Lambda}^{-1}
  =\begin{bmatrix}1/18&0\\0&1/36\end{bmatrix}.
\end{equation}

At this point, $\mat{G}$, $\mat{\Phi}$, $\mat{B}$, $\mat{A}$,
$\mat{\Lambda}$, and $\mat{W}$ are mathematically fixed.  The optimized code
does not need to materialize all of them as separate dense matrices: it stores
$\mat{B}$, the eigenvalue vector, and the already scaled forcing.

\section{Online solve for a new source}

\subsection{Step 8: receive and project the new forcing}

Suppose the new forcing on the four interior points is
\begin{equation}
  \vect{f}=\begin{bmatrix}54&9&9&18\end{bmatrix}^{\mathsf T}.
\end{equation}
Project it onto the two test functions once:
\begin{align}
  \vect{f}_M
  &=\mat{\Phi}^{\mathsf T}\vect{f}\\
  &=\begin{bmatrix}
      0.5(54)+0.5(9)+0.5(9)+0.5(18)\\
      0.5(54)-0.5(9)+0.5(9)-0.5(18)
    \end{bmatrix}\\
  &=\begin{bmatrix}45\\18\end{bmatrix}.
\end{align}

This source projection may touch the full grid, but it is performed only once
per new source, not once per latent iteration.  If the forcing has a known
affine expansion, the projections of its fixed components can also be
precomputed.

\textbf{Notation warning.} In this document, $\vect{f}_M$ means the raw
projection $\mat{\Phi}^{\mathsf T}\vect{f}$.  The implementation variable
named \texttt{f\_m} instead stores the scaled vector
$\mat{\Lambda}^{-1}\mat{\Phi}^{\mathsf T}\vect{f}$.

\subsection{Step 9: evaluate one candidate latent state}

Suppose the latent solver proposes
\begin{equation}
  \vect{z}=\begin{bmatrix}0.4&-0.7\end{bmatrix}^{\mathsf T}
\end{equation}
and the frozen latent MLP returns
\begin{equation}
  \vect{c}=\vect{h}(\vect{z})
  =\begin{bmatrix}2&-1&0.5\end{bmatrix}^{\mathsf T}.
\end{equation}

The projected left side of the PDE is
\begin{align}
  \mat{A}\vect{c}
  &=\begin{bmatrix}
      36&27&9\\0&-18&18
    \end{bmatrix}
    \begin{bmatrix}2\\-1\\0.5\end{bmatrix}\\
  &=\begin{bmatrix}
      36(2)+27(-1)+9(0.5)\\
      0(2)-18(-1)+18(0.5)
    \end{bmatrix}\\
  &=\begin{bmatrix}49.5\\27\end{bmatrix}.
\end{align}

The raw projected residual is
\begin{equation}
  \mat{A}\vect{h}(\vect{z})-\vect{f}_M
  =\begin{bmatrix}49.5\\27\end{bmatrix}
   -\begin{bmatrix}45\\18\end{bmatrix}
  =\begin{bmatrix}4.5\\9\end{bmatrix}.
\end{equation}

After applying the weights,
\begin{align}
  \vect{r}(\vect{z})
  &=\mat{W}\left(
      \mat{A}\vect{h}(\vect{z})-\vect{f}_M
    \right)\\
  &=\begin{bmatrix}4.5/18\\9/36\end{bmatrix}
   =\begin{bmatrix}0.25\\0.25\end{bmatrix}.
\end{align}
The candidate is therefore not yet a solution.

\subsection{Step 10: compute a latent update}

Automatic differentiation is needed only through the small head.  Suppose its
local Jacobian is
\begin{equation}
  \mat{J}_h
  =\frac{\partial\vect{h}}{\partial\vect{z}}
  =\begin{bmatrix}
      1&0\\
      0&1\\
      0.5&-0.5
    \end{bmatrix}
  \in\R^{3\times2}.
\end{equation}
Its three rows correspond to the $R=3$ learned-pattern coefficients, while its
two columns correspond to the $K=2$ latent variables being adjusted.
Then the residual Jacobian is
\begin{align}
  \mat{J}_r
  &=\mat{W}\mat{A}\mat{J}_h\\
  &=\begin{bmatrix}
      2.25&1.25\\
      0.25&-0.75
    \end{bmatrix}.
\end{align}

For illustration, an undamped local Newton step solves
\begin{equation}
  \mat{J}_r\,\delta\vect{z}=-\vect{r},
\end{equation}
which here gives
\begin{equation}
  \delta\vect{z}=\begin{bmatrix}-0.25\\0.25\end{bmatrix}.
\end{equation}
The production method uses Levenberg--Marquardt damping and a trust-region
limit, so it may take a smaller step.  It repeats
\begin{equation}
  \vect{z}
  \longrightarrow \vect{h}(\vect{z})
  \longrightarrow \vect{r}(\vect{z})
  \longrightarrow \delta\vect{z}
\end{equation}
until the residual or progress stopping condition is satisfied.

No spatial network, full-grid Laplacian, or $\mat{\Phi}^{\mathsf T}\mat{G}$
product is evaluated inside this loop.

\section{Convergence and final prediction}

\subsection{Step 11: reach a zero-residual coefficient vector}

Suppose the final latent state $\vect{z}_\star$ produces
\begin{equation}
  \vect{h}(\vect{z}_\star)
  =\begin{bmatrix}1&0&1\end{bmatrix}^{\mathsf T}.
\end{equation}
Then
\begin{align}
  \mat{A}\vect{h}(\vect{z}_\star)
  &=\begin{bmatrix}
      36(1)+27(0)+9(1)\\
      0(1)-18(0)+18(1)
    \end{bmatrix}\\
  &=\begin{bmatrix}45\\18\end{bmatrix}
   =\vect{f}_M.
\end{align}
Therefore
\begin{equation}
  \vect{r}(\vect{z}_\star)=\vect{0}.
\end{equation}

\subsection{Step 12: decode the full field once}

The final interior prediction is
\begin{align}
  \vect{u}_\star
  &=\mat{G}\vect{h}(\vect{z}_\star)\\
  &=
  \begin{bmatrix}
      1&0& 1\\
      0&1& 1\\
      1&1& 0\\
      2&1&-1
    \end{bmatrix}
  \begin{bmatrix}1\\0\\1\end{bmatrix}\\
  &=\begin{bmatrix}2\\1\\1\\1\end{bmatrix}.
\end{align}

Restoring the zero boundary gives the $4\times4$ field
\begin{equation}
  \mat{U}_\star
  =\begin{bmatrix}
      0&0&0&0\\
      0&2&1&0\\
      0&1&1&0\\
      0&0&0&0
    \end{bmatrix}.
\end{equation}

For completeness, the full four-point negative-Laplacian matrix is
\begin{equation}
  \mat{L}
  =9\begin{bmatrix}
      4&-1&-1&0\\
      -1&4&0&-1\\
      -1&0&4&-1\\
      0&-1&-1&4
    \end{bmatrix}.
\end{equation}
Direct multiplication verifies that this deliberately constructed toy
prediction satisfies the complete grid equation, not only the two retained
tests:
\begin{equation}
  \mat{L}\vect{u}_\star
  =\begin{bmatrix}54\\9\\9\\18\end{bmatrix}
  =\vect{f}.
\end{equation}
In a general reduced solve, a small retained test set enforces only the
selected weak equations; this full equality occurs here because the teaching
example was constructed to make it checkable end to end.

\section{What is fixed, and what changes?}

\begin{center}
\begin{tabular}{@{}>{\raggedright\arraybackslash}p{0.42\textwidth}
                    >{\raggedright\arraybackslash}p{0.48\textwidth}@{}}
\toprule
Fixed after training/precomputation & Changes for a new source or solver iterate \\
\midrule
weights of $\vect{g}$ and $\vect{h}$ & the source field $\vect{f}$ \\
spatial bank $\mat{G}$ & its projected vector $\vect{f}_M$ changes per source \\
test functions $\mat{\Phi}$ & current latent vector $\vect{z}$ \\
eigenvalues $\mat{\Lambda}$ & coefficients $\vect{h}(\vect{z})$ \\
projected matrices $\mat{B}$ and $\mat{A}$ & residual and residual Jacobian \\
weighting $\mat{W}$ & final decoded field \\
\bottomrule
\end{tabular}
\end{center}

The matrices must be rebuilt if the learned spatial bank, mesh, boundary
conditions, test space, or PDE operator changes.  A new right-hand side on the
same problem requires a new $\vect{f}_M$, but not a new $\mat{A}$.

\section{Why the implementation looks slightly shorter}

The slide form is
\begin{equation}
  \vect{r}(\vect{z})
  =\mat{W}\left[
     \mat{\Lambda}\mat{B}\vect{h}(\vect{z})-\vect{f}_M
   \right],
  \qquad
  \mat{W}=\mat{\Lambda}^{-1}.
\end{equation}
Since $\mat{W}\mat{\Lambda}=\mat{I}$, this is exactly
\begin{equation}
  \vect{r}(\vect{z})
  =\mat{B}\vect{h}(\vect{z})
   -\mat{\Lambda}^{-1}\vect{f}_M.
\end{equation}

The quadrature-free code uses this second form: it materializes
$\mat{B}=\mat{\Phi}^{\mathsf T}\mat{G}$, not the conceptual dense matrix
$\mat{A}$, and stores the forcing after inverse-eigenvalue scaling.  The
apparent disappearance of $\mat{W}$ and $\mat{\Lambda}$ is only an algebraic
fold; it does not change the residual.

\section{The Poisson method in twelve lines}

\begin{enumerate}
  \item Generate full-order Poisson training solutions.
  \item Give each training solution a trainable latent code.
  \item Train $\vect{g}$, $\vect{h}$, and the training codes to reconstruct the solutions.
  \item Freeze both neural networks.
  \item Evaluate $\vect{g}$ on the grid to construct $\mat{G}$.
  \item Construct the sine test functions $\mat{\Phi}$ and eigenvalues $\mat{\Lambda}$.
  \item Precompute $\mat{B}=\mat{\Phi}^{\mathsf T}\mat{G}$ and, conceptually,
        $\mat{A}=\mat{\Lambda}\mat{B}$.
  \item For a new source, compute $\vect{f}_M=\mat{\Phi}^{\mathsf T}\vect{f}$ once.
  \item Choose an initial latent vector $\vect{z}$.
  \item Repeatedly evaluate the small head and residual, then update $\vect{z}$.
  \item Stop at $\vect{z}_\star$ when the residual is sufficiently small.
  \item Decode once: $\vect{u}_\star=\mat{G}\vect{h}(\vect{z}_\star)$.
\end{enumerate}

\keyidea{\textbf{The central computational saving:}
the full grid is used to build the spatial bank, to project a new source, and
to return the final field.  It is not revisited during every latent-solver
iteration.  Each iteration uses only the small latent MLP and precomputed
dense matrices.}

\clearpage
\section{How this becomes a general NM-ROM}
\label{sec:general-nmrom}

The decoder used above is not only a Poisson decoder.  Its general form is
still
\begin{equation}
  \underbrace{\vect{z}}_{K\text{ latent variables}}
  \longrightarrow
  \underbrace{\vect{h}(\vect{z})}_{R\text{ coefficients}}
  \longrightarrow
  \underbrace{\mat{G}\vect{h}(\vect{z})}_{P\text{ physical-grid values}}.
\end{equation}
Here $P$ counts physical-grid values, $K$ counts the latent variables changed
by the online solver, and $R$ counts the frozen learned spatial patterns in
$\mat{G}$ and their matching coefficients in $\vect{h}(\vect{z})$.

Although $\mat{G}$ is fixed, the predicted solution is not fixed.  The
coefficient vector $\vect{h}(\vect{z})$ changes when $\vect{z}$ changes, so the
same columns of $\mat{G}$ can make many different fields.  A simple analogy is
a fixed palette of $R$ paint colors: keeping the palette fixed does not keep
the painting fixed, because the mixture changes.

There are two nested geometric objects:
\begin{itemize}
  \item All combinations $\mat{G}\vect{c}$, with any $R$-entry coefficient
        vector $\vect{c}$, form a linear space of dimension at most $R$.
  \item The neural head permits only coefficients of the form
        $\vect{c}=\vect{h}(\vect{z})$, where $\vect{z}$ has only $K$ entries.
        This usually traces a curved, $K$-dimensional surface inside the
        larger $R$-dimensional coefficient space.
\end{itemize}
That curved surface is the learned \emph{nonlinear manifold}.  The spatial
bank $\mat{G}$ can therefore be fixed while the allowed solutions depend
nonlinearly on the $K$ latent variables.

\keyidea{\textbf{What ``generalizes'' means here.}
The trained model can solve new parameter values, source fields, initial
conditions, or time states from the same problem family by finding a new
latent vector.  It does not automatically generalize to an unrelated PDE,
boundary condition, or operator.  A changed mesh requires the continuous
spatial network to be reevaluated and all projected objects to be rebuilt;
accuracy on that mesh must also be checked.}

There is an important conditional statement for linear Poisson.  If a forcing
has a finite affine expansion
\begin{equation}
  \vect{f}(\mu)=\sum_{a=1}^{A}\theta_a(\mu)\vect{f}_a,
\end{equation}
then, for a fixed invertible Poisson matrix $\mat{L}$,
\begin{equation}
  \vect{u}(\mu)
  =\sum_{a=1}^{A}\theta_a(\mu)\mat{L}^{-1}\vect{f}_a.
\end{equation}
Those exact solutions lie in the linear span of the $A$ fixed fields
$\mat{L}^{-1}\vect{f}_a$, so a nonlinear head cannot improve representation
accuracy; its possible value is reducing online cost.  The implemented
Poisson data use Gaussian sources, which are not affine in their center and
width parameters.  Therefore that theorem does not, by itself, prove that the
implemented Gaussian solution family is a low-dimensional linear subspace.
The Poisson residual precomputation remains exact because the operator is fixed
and each new source enters only through its projected right-hand side.

\section{The general offline and online recipe}

Suppose a discretized PDE is written as
\begin{equation}
  \vect{F}(\vect{u};\vect{d})=\vect{0}.
\end{equation}
The vector $\vect{u}$ has $P$ physical-grid values.  The symbol $\vect{d}$
collects the data that define the particular problem, such as a forcing,
viscosity, previous time state, or boundary data.

After auto-decoder training, freeze $\mat{G}$ and the head
$\vect{h}$.  Choose $M$ smooth test functions and store them as columns of
$\mat{\Phi}\in\mathbb{R}^{P\times M}$.  Here $M$ is the number of weak test
equations, whereas $P$ is the number of physical-grid values.  The online
solver minimizes the weighted weak residual
\begin{equation}
  \vect{r}(\vect{z};\vect{d})
  =\mat{W}\mat{\Phi}^{\mathsf T}
    \vect{F}\!\left(\mat{G}\vect{h}(\vect{z});\vect{d}\right).
\end{equation}
This says: decode a candidate field, apply the discrete PDE, project the
mismatch onto $M$ smooth test functions, and adjust the $K$ latent variables
until that projected mismatch is small.

The amount that can be precomputed is determined by the algebraic degree of
the PDE residual as a function of the state $\vect{u}$:

\begin{center}
\begin{tabular}{@{}>{\raggedright\arraybackslash}p{0.20\textwidth}
                    >{\raggedright\arraybackslash}p{0.33\textwidth}
                    >{\raggedright\arraybackslash}p{0.37\textwidth}@{}}
\toprule
State dependence & Typical term & What is stored before the online solve \\
\midrule
Degree 1: linear & diffusion, Poisson, Stokes & an $M\times R$ projected matrix \\
Degree 2: quadratic & Burgers advection, fluid convection & an $M\times R\times R$ tensor \\
Non-polynomial & $|u|$, $1/\rho$, limiters & exact polynomial pieces, plus samples or another approximation for the rest \\
\bottomrule
\end{tabular}
\end{center}

For example, a fixed linear term $\mat{L}\vect{u}$ becomes
\begin{equation}
  \mat{\Phi}^{\mathsf T}\mat{L}\mat{G}\vect{h}(\vect{z}).
\end{equation}
Everything before $\vect{h}(\vect{z})$ is one precomputed $M\times R$
matrix.  A quadratic term contains products of two entries of
$\vect{h}(\vect{z})$, so its coefficients form a three-index tensor: one
index for the $M$ test equations and two indices for pairs of the $R$ learned
patterns.

\section{Burgers: what changes from Poisson}

The implemented one-dimensional model is viscous Burgers on the interval
$[0,1]$ with zero values at both walls.  In continuous notation,
\begin{equation}
  u_t+u\,u_x=\nu u_{xx}.
\end{equation}
The viscosity $\nu$ controls diffusion.  Small $\nu$ permits a sharper moving
front; larger $\nu$ smooths it more strongly.

Backward Euler turns one time step into the nonlinear equation
\begin{equation}
  \underbrace{\vect{u}^{\,n+1}-\vect{u}^{\,n}}_{\text{change over one step}}
  +\Delta t\left[
    \underbrace{\vect{N}(\vect{u}^{\,n+1})}_{\text{advection}}
    -\underbrace{\nu\mat{\Delta}_h\vect{u}^{\,n+1}}_{\text{diffusion}}
  \right]
  =\vect{0}.
  \label{eq:burgers-step}
\end{equation}
Here $n$ labels the known old time, $n+1$ labels the unknown new time,
$\Delta t=0.005$ is the time-step size, $\nu$ is the viscosity, and
$\mat{\Delta}_h$ is the discrete second derivative.  The implementation takes
50 steps.  Therefore one initial field produces 51 saved states: the initial
state plus 50 new states.

Poisson has one solve for one source.  Burgers instead has a chain of solves.
At every time step, the previous decoded state is known and the next $K$-entry
latent vector must be found.

\subsection{The real Burgers sizes}
\label{sec:burgers-sizes}

For a concrete $N=256$ Burgers grid, two points are boundary points and
\begin{equation}
  P=256-2=254
\end{equation}
physical interior-grid values remain.  The implemented model sizes are
\begin{equation}
  K=8,\qquad R=32,\qquad M=32.
\end{equation}
Thus the online solver adjusts $K=8$ latent variables; the frozen spatial bank
contains $R=32$ learned patterns and the head returns 32 matching
coefficients; and the weak residual has $M=32$ test equations.

\begin{center}
\begin{tabular}{@{}lll@{}}
\toprule
Object & Shape at $N=256$ & Plain meaning \\
\midrule
$\vect{z}$ & $8=K$ & latent unknowns in one time-step solve \\
$\vect{h}(\vect{z})$ & $32=R$ & coefficients of the learned patterns \\
$\mat{G}$ & $254\times32=P\times R$ & learned patterns on the interior grid \\
$\mat{\Phi}$ & $254\times32=P\times M$ & 32 sine test functions on that grid \\
$\mat{B}=\mat{\Phi}^{\mathsf T}\mat{G}$ & $32\times32=M\times R$ & projected spatial bank \\
$\mat{Q}$ & $32\times32\times32=M\times R\times R$ & projected pairwise advection interactions \\
$\vect{r}$ & $32=M$ & weak residual for one candidate state \\
$\mat{J}_r$ & $32\times8=M\times K$ & residual sensitivity to the latent variables \\
\bottomrule
\end{tabular}
\end{center}

The training generator produces 512 trajectories.  Each has 51 states, so
26,112 states are available; the configuration selects at most 8,192 of them
for auto-decoder training.  Every selected state receives its own $K=8$ entry
training code.  The networks learn one $P$-by-$R$ spatial bank and a nonlinear
map from those eight latent variables to $R=32$ coefficients, exactly as in
the Poisson training description.

The Burgers spatial network starts from one coordinate $x$.  Sixty-four
learned frequencies produce 64 sine and 64 cosine values, giving 128 Fourier
features.  The path is
\begin{equation}
  \mathbb{R}^{1}
  \longrightarrow128\text{ Fourier features}
  \longrightarrow128
  \longrightarrow128
  \longrightarrow\mathbb{R}^{32}.
\end{equation}
Its 32 outputs are the $R=32$ raw learned-pattern values at $x$.  The stored
bank row is actually
\begin{equation}
  G(x,:)=s\,4x(1-x)\,\widetilde{\vect{g}}(x)^{\mathsf T},
\end{equation}
where $s$ is the fixed root-mean-square amplitude of the training data.
The factor $4x(1-x)$ makes every pattern zero at both walls.  The Burgers
latent head is
\begin{equation}
  \mathbb{R}^{8}\longrightarrow128\longrightarrow128
  \longrightarrow\mathbb{R}^{32},
\end{equation}
plus a separate linear skip from its $K=8$ latent inputs to its $R=32$
coefficient outputs.  Both hidden layers use the SiLU activation.

\section{Burgers one-time precomputation}

The $M=32$ lowest discrete sine modes are used as test functions.  They are
eigenvectors of the negative discrete Laplacian, so
\begin{equation}
  -\mat{\Delta}_h\mat{\Phi}=\mat{\Phi}\mat{\Lambda},
\end{equation}
where $\mat{\Lambda}$ contains $M$ positive eigenvalues.  First precompute
\begin{equation}
  \mat{B}=\mat{\Phi}^{\mathsf T}\mat{G}
  \in\mathbb{R}^{M\times R}.
\end{equation}
Here $M$ counts test equations and $R$ counts learned patterns.  This single
matrix handles both the state-change term and, with the eigenvalues, the
linear diffusion term in Equation~\eqref{eq:burgers-step}.

The advection term is quadratic.  At interior grid node $i$, a decoded field
value is
\begin{equation}
  u_i=\sum_{k=1}^{R}G_{ik}h_k,
\end{equation}
a backward difference gives
\begin{equation}
  (D^-u)_i=\sum_{k=1}^{R}(D^-G)_{ik}h_k.
\end{equation}
Multiplying the two expressions produces every pair of coefficients
$h_jh_k$.  Projecting that product onto test function $m$ gives
\begin{equation}
  T_{mjk}
  =\sum_{i=1}^{P}\Phi_{im}\,G_{ij}\,(D^-G)_{ik}.
  \label{eq:T-burgers}
\end{equation}
The index $m$ selects one of the $M$ test equations.  The indices $j$ and $k$
select an ordered pair of the $R$ learned spatial patterns.  Consequently
$\mat{T}$ has shape $M\times R\times R$.

The raw tensor is not symmetric because one bank factor is differenced and
the other is not.  Store the symmetrized tensor
\begin{equation}
  Q_{mjk}=T_{mjk}+T_{mkj}.
\end{equation}
For one $R$-entry coefficient vector $\vect{h}$, the complete projected
advection vector is then
\begin{equation}
  q_m(\vect{h})
  =\frac{1}{2}\vect{h}^{\mathsf T}\mat{Q}_m\vect{h},
  \qquad m=1,\ldots,M.
  \label{eq:q-burgers}
\end{equation}
There is one $R\times R$ matrix slice $\mat{Q}_m$ for each test equation.
With $M=R=32$, the production tensor contains $32^3=32{,}768$ numbers, which
occupy 256 KiB in 64-bit floating point.

\subsection{A two-coefficient tensor calculation}

This tiny arithmetic example explains the contraction in
Equation~\eqref{eq:q-burgers}.  Suppose one test equation has
\begin{equation}
  \mat{Q}_1=\begin{bmatrix}2&1\\1&4\end{bmatrix},
  \qquad
  \vect{h}=\begin{bmatrix}1\\2\end{bmatrix}.
\end{equation}
Here the teaching example has $R=2$ learned-pattern coefficients.  First,
\begin{equation}
  \mat{Q}_1\vect{h}
  =\begin{bmatrix}2&1\\1&4\end{bmatrix}
   \begin{bmatrix}1\\2\end{bmatrix}
  =\begin{bmatrix}4\\9\end{bmatrix}.
\end{equation}
Then
\begin{equation}
  q_1(\vect{h})
  =\frac12\begin{bmatrix}1&2\end{bmatrix}
    \begin{bmatrix}4\\9\end{bmatrix}
  =\frac12(22)=11.
\end{equation}
That one number is the advection contribution to one test equation.  The real
model performs the same operation for all $M=32$ test equations and all
$R=32$ coefficients at once.

\section{The Burgers residual, term by term}

Assume the old latent state $\vect{z}^{\,n}$ is known.  Its projected physical
state is computed once as
\begin{equation}
  \vect{p}^{\,n}
  =\mat{B}\vect{h}(\vect{z}^{\,n})
  =\mat{\Phi}^{\mathsf T}\vect{u}^{\,n}.
\end{equation}
The vector $\vect{p}^{\,n}$ has $M$ entries, one per test equation.  It does
not have $P$ entries because the known field has already been projected.

For a candidate new latent state $\vect{z}$, write
\begin{equation}
  \vect{c}=\vect{h}(\vect{z}).
\end{equation}
The vector $\vect{z}$ has $K=8$ latent variables and the vector $\vect{c}$ has
$R=32$ learned-pattern coefficients.  The exact small residual used by the
tensor arm is
\begin{equation}
\boxed{
  \vect{r}(\vect{z})
  =\mat{W}_{\nu}\left[
      \underbrace{\mat{B}\vect{c}-\vect{p}^{\,n}}_{\text{change from the old state}}
      +\underbrace{\Delta t\,\nu\mat{\Lambda}\mat{B}\vect{c}}_{\text{diffusion}}
      +\underbrace{\Delta t\,\vect{q}(\vect{c})}_{\text{advection}}
    \right].
}
\label{eq:burgers-small-residual}
\end{equation}
Every quantity inside the brackets has $M=32$ entries.  Reading from left to
right:
\begin{itemize}
  \item $\mat{B}\vect{c}$ is the candidate new field projected onto the test
        functions.
  \item $\vect{p}^{\,n}$ is the already projected old field.
  \item $\Delta t$ is the time-step size, here 0.005.
  \item $\nu$ is the viscosity for this trajectory.
  \item $\mat{\Lambda}\mat{B}\vect{c}$ is the projected negative-Laplacian
        contribution.
  \item $\vect{q}(\vect{c})$ is the quadratic tensor contraction from
        Equation~\eqref{eq:q-burgers}.
\end{itemize}

The fixed diagonal weighting is
\begin{equation}
  \mat{W}_{\nu}
  =\operatorname{diag}\!\left(
      \frac{1}{1+\Delta t\,\nu\lambda_1},\ldots,
      \frac{1}{1+\Delta t\,\nu\lambda_M}
    \right).
\end{equation}
It rescales high-frequency test equations, whose Laplacian eigenvalues are
larger.  It changes cheaply when the viscosity $\nu$ changes; the spatial bank
$\mat{G}$, projected bank $\mat{B}$, eigenvalues, and tensor $\mat{Q}$ remain
fixed.

For the Jacobian, the derivative of one quadratic entry is especially simple:
\begin{equation}
  \frac{\partial q_m}{\partial\vect{c}}
  =(\mat{Q}_m\vect{c})^{\mathsf T}.
\end{equation}
Stacking those rows into $\mat{D}_q(\vect{c})\in\mathbb{R}^{M\times R}$ gives
\begin{equation}
  \mat{J}_r
  =\mat{W}_{\nu}
   \left[
     \mat{B}+\Delta t\,\nu\mat{\Lambda}\mat{B}
     +\Delta t\,\mat{D}_q(\vect{c})
   \right]
   \underbrace{\frac{\partial\vect{h}}{\partial\vect{z}}}_{R\times K}.
\end{equation}
The final Jacobian has shape $M\times K$: one row per test equation and one
column per latent variable.  In the code, automatic differentiation passes
through the complete small reduced residual---the latent head plus its cached
matrix and tensor contractions.  It does not pass through a $P$-point
advection stencil.

\section{One complete Burgers rollout}

For a new initial condition and viscosity, the online calculation is:
\begin{enumerate}
  \item Receive the initial field $\vect{u}^{\,0}$ with $P$ interior values.
  \item Find an initial $K$-entry latent vector $\vect{z}^{\,0}$ whose decoded
        field $\mat{G}\vect{h}(\vect{z}^{\,0})$ fits that initial field.
  \item Compute its $M$-entry projection
        $\vect{p}^{\,0}=\mat{B}\vect{h}(\vect{z}^{\,0})$.
  \item Extrapolate the last two latent states to obtain an initial guess for
        the next state.  At the first step they are identical, so the initial
        code itself is used.
  \item Repeatedly evaluate the head, the fixed matrix products, the quadratic
        tensor contraction, and Equation~\eqref{eq:burgers-small-residual}.
        Levenberg--Marquardt adjusts only the $K=8$ latent variables.
  \item Stop when the residual is sufficiently small, progress stalls, the
        proposed step is tiny, the damping reaches its maximum after rejected
        steps, the iteration budget is exhausted, or the initial residual is
        nonfinite.  Carry the current latent state forward and record which
        stop condition occurred.
  \item Set
        $\vect{p}^{\,n+1}=\mat{B}\vect{h}(\vect{z}^{\,n+1})$ and repeat until
        all 50 steps are complete.
  \item Decode the saved latent states when the physical fields are needed.
\end{enumerate}

The inner time-step solve uses the $K$-entry latent vector, the $R$-entry head
output, the $M$-entry residual, one $M\times R$ matrix, and one
$M\times R\times R$ tensor.  It does not traverse the $P$ physical-grid values.
Fitting the initial condition and returning physical fields still touch the
grid, so the claim of grid-independent cost applies to the cached latent
solve, not to every part of the end-to-end workflow.

\section{The important upwind sign caveat}

The implemented full-order model uses sign-dependent upwinding:
\begin{equation}
  (D^{\mathrm{up}}u)_i=
  \begin{cases}
    (u_i-u_{i-1})/\Delta x, & u_i>0,\\
    (u_{i+1}-u_i)/\Delta x, & u_i\le 0.
  \end{cases}
\end{equation}
The advection at node $i$ is $u_i(D^{\mathrm{up}}u)_i$.  A fixed
backward-difference tensor is guaranteed to represent this expression exactly
for a nonnegative decoded field.  At $u_i=0$, the forward/backward choice is
irrelevant because the product is zero.  Once negative values occur, the
stencil changes and the tensor generally differs from the full sign-upwind
operator, although special cancellations can still occur.

The same operator can be written as
\begin{equation}
  \vect{N}_{\mathrm{upwind}}(\vect{u})
  =\vect{u}\odot D^c\vect{u}
   -\frac{\Delta x}{2}|\vect{u}|\odot\mat{\Delta}_h\vect{u}.
\end{equation}
The first term is quadratic and can be stored in a tensor.  The second contains
$|\vect{u}|$, which is not a polynomial in the state, so it still needs
sampling or another approximation for genuinely sign-changing fields.

This distinction leads to three logically different Burgers models:
\begin{itemize}
  \item \textbf{Sign-upwind plus a plain backward tensor:} guaranteed exact on
        the nonnegative part of state space; its general mismatch on
        sign-changing states must be measured, not hidden.
  \item \textbf{Sign-upwind in split form:} precompute the quadratic central
        term and sample only the $|u|$ stabilizer.
  \item \textbf{A skew-symmetric or centered quadratic discretization:}
        representable exactly by a tensor for either sign, but it is a
        different full-order discretization and must be judged against its own
        full-order solution.
\end{itemize}

For the current positive-blob family, the full-order fields are intended to
remain nonnegative, while decoded candidates can undershoot slightly below
zero.  Therefore the plain tensor is algebraically exact on nonnegative
decoded fields and only near the full sign-upwind oracle when those
undershoots occur.  It should not be described as universally exact for
upwind Burgers.

\section{Poisson and Burgers side by side}

\begin{center}
\begin{tabular}{@{}>{\raggedright\arraybackslash}p{0.20\textwidth}
                    >{\raggedright\arraybackslash}p{0.34\textwidth}
                    >{\raggedright\arraybackslash}p{0.34\textwidth}@{}}
\toprule
Question & Poisson & Burgers \\
\midrule
Unknown online & one $K$-entry latent vector for a new source & one $K$-entry latent vector at each time step \\
Known problem data & projected source & viscosity and projected previous state \\
State degree & linear, degree 1 & linear diffusion plus quadratic advection \\
Main cached object & $M\times R$ matrix & $M\times R$ matrix and $M\times R\times R$ tensor \\
Full grid in the inner solve? & no & no for the tensor arm \\
Main limitation & new operator or boundary data require rebuilding & non-polynomial sign-upwind part is not one global tensor \\
\bottomrule
\end{tabular}
\end{center}

\keyidea{\textbf{Answer to ``Is a fixed matrix generalizable?''}
Yes, within the problem family for which it was built.  The fixed matrix is
not a stored answer; it is a compact representation of how the fixed PDE
operator acts on the learned spatial patterns.  New data change the latent
coefficients and therefore change the solution.  If the operator has affine
parameter dependence, one matrix can be stored for each fixed component and
combined online.  A new PDE, boundary condition, non-affine operator, or
unvalidated state family requires new projected objects and possibly new
training.}

\section{The full NM-ROM pipeline in fifteen lines}

\begin{enumerate}
  \item Generate full-order solution snapshots from the intended PDE family.
  \item Give each selected snapshot a trainable $K$-entry latent code.
  \item Train the spatial network, latent head, and codes to reconstruct the snapshots.
  \item Freeze both networks.
  \item Evaluate the spatial network to build the $P\times R$ bank $\mat{G}$.
  \item Choose $M$ smooth test functions and build $\mat{\Phi}$.
  \item Precompute an $M\times R$ matrix for every linear state term.
  \item Precompute an $M\times R\times R$ tensor for every quadratic state term.
  \item Identify every non-polynomial term that still requires sampling.
  \item Receive new problem data: a source for Poisson, or an initial condition and viscosity for Burgers.
  \item Obtain an initial $K$-entry latent guess.
  \item Evaluate only the small head and cached reduced operators inside the latent solve.
  \item Use the $M\times K$ residual Jacobian to update the latent vector.
  \item For an unsteady PDE, carry the accepted latent state into the next time step.
  \item Decode $\mat{G}\vect{h}(\vect{z})$ when the $P$-value physical field is needed.
\end{enumerate}

\clearpage
\section*{Plain-language glossary}
\addcontentsline{toc}{section}{Plain-language glossary}

\begingroup
\setlength{\parskip}{0.5em}

\textbf{$K$ --- latent variables.}
The number of entries in $\vect{z}$ that the online solver adjusts.

\textbf{$R$ --- learned spatial patterns.}
The number of columns in $\mat{G}$ and, equivalently, the number of
coefficients returned by $\vect{h}(\vect{z})$.

\textbf{$M$ --- test equations.}
The number of weak test equations retained by the ROM.

\textbf{$P$ --- physical-grid values.}
The number of interior grid values stored in one solution vector $\vect{u}$.

\textbf{$\vect{g}(\vect{x})$ --- spatial neural network.}
Returns the raw learned-pattern values at coordinate $\vect{x}$.

\textbf{$\vect{h}(\vect{z})$ --- latent neural network.}
Returns the coefficients that combine the learned spatial patterns.

\textbf{$\mat{G}$ --- frozen spatial bank.}
Contains the learned patterns evaluated at every interior grid point,
including the fixed output scale and hard boundary mask.

\textbf{$\mat{\Phi}$ --- test-function matrix.}
Its columns are the selected sine test functions.

\textbf{$\mat{B}$ --- projected bank.}
This is $\mat{\Phi}^{\mathsf T}\mat{G}$.

\textbf{$\mat{\Lambda}$ --- eigenvalue matrix.}
The diagonal matrix of Laplacian eigenvalues for the sine test modes.

\textbf{$\mat{A}$ --- projected Poisson operator.}
This equals $\mat{\Lambda}\mat{B}$ in this setting.

\textbf{$\mat{W}$ --- residual weighting.}
The fixed inverse-eigenvalue scaling.

\textbf{$\vect{f}_M$ --- projected forcing.}
The new forcing projected onto the retained test functions.

\textbf{$\vect{r}(\vect{z})$ --- residual.}
The weighted mismatch between the decoded field's projected Poisson response
and the projected forcing, or between two consecutive states and the projected
time-dependent PDE terms for Burgers.

\textbf{$\mat{T}$ --- raw Burgers interaction tensor.}
Stores the projection of every ordered pair consisting of one learned pattern
and one backward-differenced learned pattern.

\textbf{$\mat{Q}$ --- symmetric Burgers interaction tensor.}
The raw tensor plus its transpose in the two learned-pattern indices.  It is
used in the quadratic contraction and its derivative.

\textbf{$\vect{q}(\vect{h})$ --- projected advection.}
The $M$-entry vector produced by contracting the Burgers tensor with the
$R$-entry coefficient vector twice.

\textbf{$\Delta t$ --- time-step size.}
The amount of simulated time advanced by one backward-Euler solve.

\textbf{$\nu$ --- viscosity.}
The Burgers parameter controlling the strength of diffusion.

\textbf{FOM --- full-order model.}
The reference discretized PDE solver that evolves all physical-grid values.

\textbf{LM --- Levenberg--Marquardt.}
The damped nonlinear least-squares method used to adjust the latent variables.

\textbf{NM-ROM --- nonlinear-manifold reduced-order model.}
A reduced solver whose admissible states are selected by a nonlinear map from
the latent vector, here $\vect{u}=\mat{G}\vect{h}(\vect{z})$.

\textbf{Polynomial degree in the state.}
The highest number of state factors multiplied together in a residual term:
one for a linear term, two for a quadratic term.

\textbf{Sign-upwind.}
An advection difference that switches between a backward and forward stencil
according to the sign of the local velocity.

\textbf{Tensor contraction.}
Summing the stored pairwise interaction coefficients against every product
$h_jh_k$ to produce the projected quadratic term.

\textbf{Auto-decoder.}
A decoder trained with one free latent code per training snapshot, without a
learned encoder.

\textbf{Offline.}
Training and precomputation performed before solving new source instances.

\textbf{Online.}
The latent solve and final field reconstruction for a new source.

\endgroup

\end{document}
